% !TeX root = ../../tesis.tex

% =============================================================================
%  3.3.0 — Introducción a los Indicadores de Análisis
% =============================================================================

El análisis computacional de la red de transporte público de la \cdmx{} y su
zona metropolitana se articula a través de ocho indicadores cuantitativos
organizados en cuatro fases de construcción progresiva. Esta jerarquía no es
arbitraria: cada indicador depende de los insumos producidos por los anteriores,
de modo que la arquitectura metodológica constituye en sí misma un modelo de
grafo que crece en complejidad desde el análisis espacial puro hasta las pruebas
de estrés topológico. A continuación se presenta cada uno de los indicadores con
su fundamento matemático, la fuente académica que lo sustenta, y un ejemplo
práctico calibrado con datos del área de estudio.

Los ocho indicadores están distribuidos entre los Objetivos Específicos 1 al 5
(§\ref{subsec:obj-especificos}): los indicadores de la Fase 1 ---$C$, $C_i$, $DI$ y $CF_h$--- alimentan los
Objetivos Específicos 2 y 4; el de la Fase 2 ---$W_{s,h}$--- refina el grafo
del Objetivo Específico 1; el Tiempo de Viaje
Promedio $T$ y la Centralidad de Intermediación $B(v)$ responden a los Objetivos
Específicos 2 y 3 respectivamente; y la Robustez Geométrica $\Delta E$ integra
los Objetivos Específicos 3 y 5 al cuantificar el efecto de la propuesta anillar
sobre la vulnerabilidad estructural de la red. La tabla de síntesis al final de
esta sección (§\ref{subsec:resumen-indicadores}) concentra estas
correspondencias de manera visual.

\bigskip
\noindent\textbf{Nota aclaratoria sobre el alcance de los indicadores.}
Todos los indicadores enumerados en este capítulo se definen en su fundamento
conceptual y matemático. No obstante, el indicador de
\textbf{Robustez Geométrica $\Delta E$} constituye una excepción en cuanto a su
demostración aplicada: por los tiempos del presente estudio y el enfoque de los
resultados esperados en las capas de código, $\Delta E$ no se implementa ni
valida empíricamente en los capítulos de desarrollo tecnológico. Su definición
formal se preserva aquí como marco teórico; la implementación computacional y la
validación sobre la red de la \zmvm{} se proponen como línea prioritaria de
investigación futura.
